\chapter{TINJAUAN PUSTAKA}

\section{Penelitian Terdahulu}
Beberapa penelitian sebelumnya telah mengeksplorasi permasalahan serupa. \citet{riedmiller1994} mengusulkan algoritma \textit{gradient descent} untuk meminimumkan \textit{mean square error} (MSE) pada pelatihan jaringan saraf tiruan. Pendekatan ini kemudian dikembangkan lebih lanjut oleh beberapa peneliti lain \citep{calvert2005,aaron2014}.

Sementara itu, meskipun dianggap sebagai pendekatan yang relatif baru, algoritma \textit{artificial bee colony} (ABC) telah digunakan secara luas dalam pencarian penyelesaian optimum \citep{shukran2011}.

\subsection{Ringkasan perbandingan penelitian terdahulu}
Tabel~\ref{tab:perbandingan} merangkum perbandingan pendekatan, dataset, dan metrik evaluasi yang digunakan pada beberapa penelitian terdahulu yang relevan dengan penelitian ini.

\begin{table}[H]
\centering
\caption{Perbandingan penelitian terdahulu}
\label{tab:perbandingan}
\begin{tabular}{@{}lllc@{}}
\toprule
Peneliti & Metode & Dataset & Metrik \\
\midrule
\citet{riedmiller1994} & Gradient Descent & Sintetis & MSE \\
\citet{shukran2011}    & Artificial Bee Colony & Benchmark & Akurasi \\
\citet{russell2020}    & Pembelajaran Mesin & Umum & F1-score \\
\bottomrule
\end{tabular}
\end{table}

\section{Landasan Teori}

\subsection{Data science dan kecerdasan buatan}
Data science merupakan bidang interdisipliner yang menggabungkan statistik, ilmu komputer, dan pengetahuan domain untuk mengekstraksi wawasan dari data. Kecerdasan buatan (\textit{artificial intelligence}) adalah cabang ilmu komputer yang berfokus pada pembangunan sistem yang mampu meniru kemampuan kognitif manusia \citep{russell2020}.

\subsection{Metode yang relevan dengan penelitian ini}
Uraikan di sini teori dan konsep matematis dari metode/algoritma yang menjadi dasar pendekatan yang diusulkan pada Bab 3, misalnya arsitektur model pembelajaran mesin, fungsi \textit{loss}, atau kerangka evaluasi yang digunakan.
